\documentclass[11pt]{article}
\usepackage[margin=1in]{geometry}
\usepackage{amsmath,amssymb,amsthm,graphicx}
\usepackage[hidelinks]{hyperref}
\newtheorem{theorem}{Theorem}
\newtheorem{lemma}{Lemma}
\newcommand{\R}{\mathbb R}
\newcommand{\one}{\mathbf 1}
\newcommand{\WMP}{\mathrm{WMP}^{+}}
\setlength{\parindent}{0pt}
\setlength{\parskip}{0.65em}

\title{The Positive Schur Property Does Not Characterize Universal $\WMP$ Codomains}
\author{Run \texttt{fa\_banach\_001}, lane 18}
\date{11 August 2026}

\begin{document}
\maketitle

\textbf{Status.} Candidate counterexample, likely valid, for human review.

\section{Source question}

Luiz and Miranda introduce the positive weak maximizing property: a pair of
Banach lattices $(E,F)$ has $\WMP$ when every positive operator $T:E\to F$
with a positive non-weakly-null maximizing sequence is norm-attaining. After
recalling that a Banach space $Y$ has the Schur property exactly when $(X,Y)$
has the ordinary WMP for every reflexive $X$, they ask whether there is a
similar characterization for $\WMP$ using the positive Schur property
\cite[Section~3, PDF p.~12]{LuizMiranda}.

\begin{center}
\includegraphics[width=\textwidth]{figures/open_problem_crop.png}
\end{center}

The direct candidate characterization is:
\begin{quote}
$F$ has the positive Schur property if and only if $(E,F)$ has $\WMP$ for
every reflexive Banach lattice $E$.
\end{quote}
The source proves the proposed sufficiency only when the unit ball of $E$ is
sequentially absolutely weakly compact, a condition stronger than
reflexivity, and proves the relevant necessity result for order-continuous
codomains. The construction below disproves sufficiency for reflexive domains.

\section{Claimed contribution}

\begin{theorem}\label{thm:main}
For every $\delta>0$ there is an order-continuous lattice norm
$\|\!|\cdot\|\!|$ on $L_1[0,1]$ such that
\[
 \|h\|_1\leq \|\!|h\|\!|\leq(1+\delta)\|h\|_1
 \qquad(h\in L_1[0,1]),
\]
and the resulting Banach lattice $F=(L_1[0,1],\|\!|\cdot\|\!|)$ has the
positive Schur property, but
\[
 (L_2[0,1],F)
\]
fails $\WMP$. Consequently the positive Schur property does not characterize
the codomains $F$ for which $(E,F)$ has $\WMP$ for every reflexive Banach
lattice $E$.
\end{theorem}

Thus the failure persists under arbitrarily small equivalent lattice
renormings of the standard $L_1$ norm.

\section{Proof intuition}

An equivalent lattice renorming preserves the positive Schur property, but it
can encode a moving family of positive functionals whose norms increase to an
unattained limit. Rademacher half-sets provide exactly such a family. Their
oscillatory parts vanish weakly against each fixed $L_2$ vector, preventing
any single vector from attaining the limiting operator norm. On the other
hand, each density nearly norms itself, so its normalization gives a
maximizing sequence. Every density also has a fixed positive constant part;
therefore that maximizing sequence converges weakly to a nonzero constant and
is not weakly null. This is precisely the witness that violates $\WMP$.

\section{Construction and proof}

Work over the real scalars. Let $(r_n)$ be the Rademacher functions on
$[0,1]$, normalized so that $r_n$ takes the values $-1$ and $1$ on sets of
measure $1/2$. Put
\[
 A_n=\{t:r_n(t)=1\},\qquad \one_{A_n}=\frac{\one+r_n}{2}.
\]
Fix $c$ with $0<c\leq\delta$, choose any strictly increasing sequence
$0<c_n<c$ with $c_n\to c$ (for instance $c_n=cn/(n+1)$), and define
\[
 g_n=\one+c_n\one_{A_n}.
\]
For $h\in L_1[0,1]$, set
\begin{equation}\label{eq:Fnorm}
 \|\!|h\|\!|=\sup_{n\geq1}\int_0^1 g_n|h|.
\end{equation}

\begin{lemma}\label{lem:lattice}
Equation~\eqref{eq:Fnorm} defines an order-continuous Banach lattice norm on
$L_1[0,1]$ satisfying
\[
 \|h\|_1\leq\|\!|h\|\!|\leq(1+c)\|h\|_1.
\]
The resulting Banach lattice $F$ has the positive Schur property.
\end{lemma}

\begin{proof}
Since $\one\leq g_n\leq(1+c)\one$, the displayed norm equivalence is
immediate. Positivity of the $g_n$, the inequality
$|h+k|\leq|h|+|k|$, and monotonicity of the integral prove the norm and
lattice-norm axioms. Completeness follows from equivalence with the ordinary
$L_1$ norm. Since the ordinary $L_1$ norm is order continuous, the same
equivalence shows that $\|\!|\cdot\|\!|$ is order continuous as well.

Equivalent norms have the same continuous linear functionals and hence the
same weak topology. If $(h_j)\subset F_+$ is weakly null, then the constant
functional $\one\in L_\infty$ gives
\[
 \|h_j\|_1=\int_0^1h_j\longrightarrow0.
\]
Consequently
$\|\!|h_j\|\!|\leq(1+c)\|h_j\|_1\to0$. This is the positive Schur
property.
\end{proof}

Consider the natural inclusion
\[
 T:L_2[0,1]\longrightarrow F,\qquad Tf=f.
\]
It is positive. Write
\[
 d_n=\|g_n\|_2
 =\left(1+c_n+\frac{c_n^2}{2}\right)^{1/2},
 \qquad
 L=\left(1+c+\frac{c^2}{2}\right)^{1/2}.
\]
Then $d_n\to L$ and $d_n<L$ for every $n$.

\begin{lemma}\label{lem:operator}
The operator $T$ has norm $L$ and does not attain its norm.
\end{lemma}

\begin{proof}
For $f\in L_2$, Cauchy--Schwarz gives
\[
 \|\!|Tf\|\!|
 =\sup_n\int_0^1g_n|f|
 \leq\sup_n\|g_n\|_2\|f\|_2
 =L\|f\|_2.
\]
For the reverse inequality let $x_n=g_n/d_n$. Then $x_n\geq0$,
$\|x_n\|_2=1$, and
\[
 \|\!|Tx_n\|\!|
 \geq\int_0^1g_nx_n=d_n\longrightarrow L.
\]
Thus $\|T\|=L$.

Fix now any $f\in L_2$ with $\|f\|_2=1$ and put
$a_n=\int g_n|f|$. Since
\[
 g_n=\left(1+\frac{c_n}{2}\right)\one+\frac{c_n}{2}r_n
\]
and $(r_n)$ is an orthonormal sequence in $L_2$, its Fourier coefficients
against the fixed vector $|f|\in L_2$ tend to zero. Hence
\begin{equation}\label{eq:fixedlimit}
 a_n\longrightarrow
 \left(1+\frac c2\right)\|f\|_1
 \leq1+\frac c2.
\end{equation}
But
\[
 L^2-\left(1+\frac c2\right)^2=\frac{c^2}{4}>0,
\]
so the limit in~\eqref{eq:fixedlimit} is strictly less than $L$. Also,
for each fixed $n$,
\[
 a_n\leq\|g_n\|_2\|f\|_2=d_n<L.
\]
Because $(a_n)$ converges to a value below $L$, its tail is uniformly below
$L$; the finitely many remaining terms are each below $L$. Therefore
\[
 \|\!|Tf\|\!|=\sup_na_n<L.
\]
No unit vector $f$ attains $\|T\|=L$.
\end{proof}

It remains to verify that the already constructed sequence is the witness
required by $\WMP$. From
\[
 x_n=\frac{1+c_n/2}{d_n}\one+\frac{c_n}{2d_n}r_n
\]
and weak nullity of the orthonormal sequence $(r_n)$ in $L_2$, we obtain
\[
 x_n\ \xrightarrow{\,w\,}\ \frac{1+c/2}{L}\one\neq0.
\]
Thus $(x_n)$ is positive and non-weakly-null. The estimates in the proof of
Lemma~\ref{lem:operator} show
\[
 d_n\leq\|\!|Tx_n\|\!|\leq\|T\|=L,
\]
so $(x_n)$ is a maximizing sequence for $T$. The positive operator $T$ is
not norm-attaining but admits a positive non-weakly-null maximizing sequence.
Hence $(L_2,F)$ fails $\WMP$, proving Theorem~\ref{thm:main}.

\section{Verification and limitations}

The separate adversarial report \texttt{verification.md} audits the Banach
lattice norm, the positive Schur property, the exact operator norm, the strict
nonattainment argument, and the maximizing sequence. No numerical or symbolic
computation is used as proof. The only analytic ingredients are
Cauchy--Schwarz, weak nullity of an orthonormal sequence, and elementary norm
equivalence.

The construction answers the direct positive-Schur characterization
negatively. It does not give an intrinsic replacement characterization of
those $F$ for which $(E,F)$ has $\WMP$ for every reflexive Banach lattice
$E$.

\section{Novelty check}

The bounded search on 11 August 2026 covered the run's registry, solution,
attempt, and proof-gap indexes; exact searches for arXiv:2409.14625 and its
title; the phrases ``positive weak maximizing property'' and ``positive Schur
property''; the published JMAA article; and arXiv searches for the exact
question and close variants. The published 2026 version still presents the
characterization as a natural question. No later proof or counterexample was
found. This supports, but cannot certify, novelty.

\begin{thebibliography}{9}
\bibitem{LuizMiranda}
J. L. P. Luiz and V. C. C. Miranda,
\emph{On norm-attaining positive operators between Banach lattices},
J. Math. Anal. Appl. 558 (2026), Paper No. 130379;
arXiv:2409.14625,
\url{https://arxiv.org/abs/2409.14625}.
\end{thebibliography}

\end{document}
